\documentclass[10pt,a4paper]{article}
\usepackage{fontspec}
\defaultfontfeatures{Ligatures=TeX}
\setmainfont{PT Serif}
\setsansfont{PT Sans}
\setmonofont{Menlo}[Scale=0.85]
\usepackage{polyglossia}
\setdefaultlanguage{russian}
\newfontfamily\cyrillicfont{PT Serif}[Script=Cyrillic]
\newfontfamily\cyrillicfontsf{PT Sans}[Script=Cyrillic]
\newfontfamily\cyrillicfonttt{Menlo}[Scale=0.85]
\usepackage[a4paper,top=15mm,bottom=16mm,left=15mm,right=15mm]{geometry}
\usepackage{graphicx,amsmath,caption,needspace,booktabs,tabularx,array,fancyvrb}
\PassOptionsToPackage{obeyspaces}{url}
\usepackage[hidelinks]{hyperref}
\usepackage{xurl,microtype}
\graphicspath{{figures/}}
\hypersetup{pdftitle={Технология MapReduce. Вариант 2},pdfauthor={Бондаренко Анна Андреевна}}
\setlength{\parindent}{0pt}
\setlength{\parskip}{3pt}
\setlength{\emergencystretch}{2em}
\linespread{1.06}
\captionsetup{font=small,labelfont=normalfont,labelsep=colon,justification=centering,singlelinecheck=false,skip=3pt}
\renewcommand{\figurename}{Рисунок}
\pagestyle{plain}
\newcolumntype{Y}{>{\raggedright\arraybackslash}X}
\newcolumntype{L}[1]{>{\raggedright\arraybackslash}p{#1}}
\newcommand{\heading}[1]{\Needspace{4\baselineskip}{\LARGE\sffamily\bfseries #1\par}\vspace{1mm}\hrule\vspace{2mm}}
\newcommand{\block}[1]{\Needspace{3\baselineskip}\vspace{1.5mm}{\sffamily\bfseries #1\par}\vspace{0.3mm}}
\newcommand{\shot}[2]{\par\vspace{2pt}\begin{minipage}{\linewidth}\centering\includegraphics[width=\linewidth]{#1.png}\captionof{figure}{#2}\end{minipage}\par\vspace{3pt}}
\newcommand{\file}[1]{\nolinkurl{#1}}
\begin{document}
\renewcommand{\figurename}{Рисунок}
\begin{titlepage}
\centering
\vspace*{15mm}
{\Large\sffamily\bfseries Университет ИТМО\par}
\vspace{35mm}
{\LARGE\sffamily\bfseries Отчёт по домашнему заданию № 3\par}
\vspace{8mm}
{\Large Технология MapReduce\par}
\vspace{12mm}
{\large Дисциплина «Инженерия данных»\par}
\vspace{12mm}
{\large\bfseries Вариант № 2\par}
\vspace{3mm}
{\large Программист, разработчик на Python\par}
\vfill
\begin{flushright}
{\large Выполнила:\\[2mm]Бондаренко Анна Андреевна}
\end{flushright}
\vspace*{35mm}
\end{titlepage}
\setcounter{page}{2}


\heading{1. RustFS и исходные данные}

Использован кластер RustFS из предыдущего ДЗ: четыре контейнера \file{rustfs-lab-node1}--\file{rustfs-lab-node4}. S3 API доступен на порту \file{19000}, консоль --- на \file{19001}. Конфигурация сохранена в \file{compose.yaml}.

\shot{01_containers_started}{Запуск четырёх контейнеров RustFS}


В бакет \file{hh-lake} загружен предоставленный файл \file{it_vacancies_2024.csv}. Размер --- 847\,619\,384 байта, или 808,35 MiB. Путь объекта: \file{s3://hh-lake/raw/it_vacancies_2024.csv}. После загрузки объект полностью прочитан из S3. Размер и SHA-256 совпали с исходником.

\shot{02_raw_object}{Исходный CSV в бакете hh-lake}


Из 229\,388 записей выбраны вакансии с ролью «Программист, разработчик» и отдельным словом Python в названии без учёта регистра. Python только в списке навыков недостаточно для отбора.

\begingroup\small\setlength{\tabcolsep}{3pt}\renewcommand{\arraystretch}{1.18}
\begin{tabularx}{\linewidth}{Yr}
\toprule
Показатель & Значение \\
\midrule
Вакансий варианта № 2 & 3\,892 \\
С указанными навыками & 2\,984 \\
Без указанных навыков & 908 \\
С любой указанной зарплатой & 1\,446 \\
С рублёвой зарплатой на руки & 912 \\
\bottomrule
\end{tabularx}
\endgroup


\newpage\heading{2. Подсчёт MapReduce}

\begin{Verbatim}[fontsize=\footnotesize]
python scripts/run_mapreduce.py --upload
\end{Verbatim}


CSV читается непосредственно из S3 чанками по 1\,000 записей. Работают два процесса-маппера, всего обработано 230 чанков. Навыки приводятся к нижнему регистру, пробелы по краям и повторы внутри вакансии удаляются. Мапперы считают частоты сочетаний, родительский процесс суммирует их по ключам.

Вакансия с $n$ уникальными навыками даёт $n$ отдельных навыков, $\binom{n}{2}$ пар и $\binom{n}{3}$ троек. Например, 10 навыков дают 45 пар и 120 троек. Поэтому число сочетаний растёт быстрее числа навыков.

\shot{03_mapreduce_run}{Запуск MapReduce с чтением S3}


Результаты сохранены в \file{hh-lake/results/mapreduce}. Для размеров сочетаний 1, 2 и 3 манифест содержит соответственно 1, 3 и 15 gzip-порций. Их порядок и контрольные суммы указаны в \file{manifest.json}. Порции ограничены 60 KiB из-за ошибки записи крупных объектов в используемой сборке RustFS на Docker Desktop.

\shot{12_s3_results}{Результаты для навыков, пар и троек в RustFS}


\newpage\heading{3. DuckDB и сверка}

\begin{Verbatim}[fontsize=\footnotesize]
python scripts/run_duckdb.py
python scripts/compare_results.py
\end{Verbatim}


DuckDB независимо читает тот же объект S3. Условие варианта в \file{sql/prepare.sql}:

\begin{Verbatim}[fontsize=\footnotesize]
where trim(role) = 'Программист, разработчик'
  and regexp_matches(
      lower(name),
      '(^|[^\p{L}\p{N}_])python([^\p{L}\p{N}_]|$)'
  )
\end{Verbatim}


Для пар и троек навыки соединяются по \file{vacancy_id} с условиями \file{a.skill < b.skill} и \file{b.skill < c.skill}, чтобы исключить перестановки и повторы.

\shot{04_duckdb_run}{Независимый расчёт DuckDB}


\begingroup\small\setlength{\tabcolsep}{3pt}\renewcommand{\arraystretch}{1.18}
\begin{tabularx}{\linewidth}{cYYYY}
\toprule
$k$ & Ключей MR & Ключей SQL & Вхождений MR & Вхождений SQL \\
\midrule
1 & 1\,269 & 1\,269 & 19\,727 & 19\,727 \\
2 & 17\,798 & 17\,798 & 81\,302 & 81\,302 \\
3 & 115\,106 & 115\,106 & 291\,454 & 291\,454 \\
\bottomrule
\end{tabularx}
\endgroup


Вхождения --- сумма частот всех сочетаний соответствующего размера. Проверены все ключи, частоты и ID выбранных вакансий. Расхождений нет.

\shot{05_comparison}{Полная сверка MapReduce и DuckDB}


\block{Повтор маппера}

\begin{Verbatim}[fontsize=\footnotesize]
python scripts/run_mapreduce.py --inject-failure --output mapper_failure
python scripts/compare_results.py --left mapreduce --right mapper_failure
\end{Verbatim}


На чанке 0 смоделирована ошибка задачи до передачи результата редьюсеру. Чанк повторён один раз. Результат совпал с обычным запуском, двойного учёта нет.

\shot{06_mapper_failure}{Ошибка маппера и повтор чанка}


\shot{07_failure_comparison}{Сверка после повтора}


\newpage\heading{4. Время и память}

\begin{Verbatim}[fontsize=\footnotesize]
python scripts/benchmark.py --runs 3
\end{Verbatim}


Каждый подход запущен трижды в отдельном процессе. Измерено время запуска Python, чтения S3, расчёта и записи локальных результатов. Загрузка в S3 и внешняя сверка в замер не входят. Кеши не сбрасывались, \file{httpfs} установлен заранее.

Память --- максимальная наблюдавшаяся сумма RSS родителя и дочерних процессов с паузой опроса 50 мс. Общие страницы могут учитываться повторно. Память RustFS и Docker VM не измеряется.

\begingroup\small\setlength{\tabcolsep}{3pt}\renewcommand{\arraystretch}{1.18}
\begin{tabularx}{\linewidth}{cYYYY}
\toprule
Попытка & MR, с & MR, MiB & DuckDB, с & DuckDB, MiB \\
\midrule
1 & 21,154 & 364,89 & 18,839 & 1\,061,55 \\
2 & 19,545 & 369,41 & 15,814 & 1\,246,50 \\
3 & 20,616 & 338,70 & 14,870 & 1\,309,12 \\
\bottomrule
\end{tabularx}
\endgroup


\begingroup\small\setlength{\tabcolsep}{3pt}\renewcommand{\arraystretch}{1.18}
\begin{tabularx}{\linewidth}{Yrr}
\toprule
Подход & Медиана времени, с & Максимум RSS, MiB \\
\midrule
MapReduce & 20,616 & 369,41 \\
DuckDB & 15,814 & 1\,309,12 \\
\bottomrule
\end{tabularx}
\endgroup


Во всех шести попытках результаты совпали. Медианное время DuckDB меньше на 23,3\%, максимальная сумма RSS больше в 3,54 раза. DuckDB использует буферы SQL-движка и группировок. MapReduce разбирает CSV и суммирует частоты в Python, а также передаёт чанки между процессами. В этих запусках выигрыш по времени сопровождается большим расходом памяти.

Одиночный запуск MapReduce с \file{--upload} включает загрузку результатов, поэтому его время отличается от таблицы профилирования.

\shot{08_benchmark}{Три запуска каждого подхода с измерением времени и памяти}


\newpage\heading{5. Навыки и зарплаты}

Анализ выполнен на той же выборке из 3\,892 Python-вакансий. Использованы 912 зарплат в рублях с \file{gross=false}. Для вилки берётся середина, при одной границе --- указанная сумма. Валюты и зарплаты до вычета налогов с этой выборкой не смешиваются.

Для пары база --- более частый навык, для тройки --- более частая входящая пара. Частота определяется по всей выборке, медианы --- по зарплатной части. База включает целевую группу. В целевой группе не менее 10 зарплат.

Разница рассчитывается как $(M_{\text{сочетания}}/M_{\text{базы}}-1)\cdot100\%$.

\block{5.1. Пары навыков}

\begin{Verbatim}[fontsize=\footnotesize]
python scripts/analysis.py salary --size 2
\end{Verbatim}


\begingroup\small\setlength{\tabcolsep}{3pt}\renewcommand{\arraystretch}{1.18}
\begin{tabularx}{\linewidth}{YYrrcr}
\toprule
Сочетание & База & \shortstack[r]{Медиана\\базы, руб.} & \shortstack[r]{Медиана\\группы, руб.} & $n$ & Изм., \% \\
\midrule
bottle + python & python & 150\,000 & 375\,000 & 11 & +150,0 \\
python + zeromq & python & 150\,000 & 375\,000 & 11 & +150,0 \\
asyncio + redis & redis & 175\,000 & 400\,000 & 19 & +128,6 \\
git + html & git & 140\,000 & 57\,500 & 20 & -58,9 \\
css + git & git & 140\,000 & 61\,250 & 14 & -56,2 \\
java + python & python & 150\,000 & 67\,500 & 12 & -55,0 \\
\bottomrule
\end{tabularx}
\endgroup


$n$ --- число зарплат в целевой группе. В базах Python, Redis и Git соответственно 709, 149 и 275 зарплат.

\shot{09_salary_pairs}{Пары навыков с положительной и отрицательной разницей медиан}


\newpage\heading{5. Навыки и зарплаты}

\block{5.2. Тройки навыков}

\begin{Verbatim}[fontsize=\footnotesize]
python scripts/analysis.py salary --size 3
\end{Verbatim}


\begingroup\small\setlength{\tabcolsep}{3pt}\renewcommand{\arraystretch}{1.18}
\begin{tabularx}{\linewidth}{YYrrcr}
\toprule
Сочетание & База & \shortstack[r]{Медиана\\базы, руб.} & \shortstack[r]{Медиана\\группы, руб.} & $n$ & Изм., \% \\
\midrule
elasticsearch + python + redis & python + redis & 180\,000 & 420\,000 & 11 & +133,3 \\
kafka + python + redis & python + redis & 180\,000 & 420\,000 & 11 & +133,3 \\
asyncio + postgresql + redis & postgresql + redis & 180\,000 & 400\,000 & 18 & +122,2 \\
git + html + python & git + python & 135\,000 & 57\,500 & 20 & -57,4 \\
css + git + python & git + python & 135\,000 & 61\,250 & 14 & -54,6 \\
css + python + sql & python + sql & 150\,000 & 70\,000 & 15 & -53,3 \\
\bottomrule
\end{tabularx}
\endgroup


$n$ --- число зарплат в целевой группе. В базах Python + Redis, PostgreSQL + Redis, Git + Python и Python + SQL соответственно 138, 104, 260 и 222 зарплаты.

\shot{10_salary_triples}{Тройки навыков с положительной и отрицательной разницей медиан}


Положительная разница связана с добавлением Bottle, ZeroMQ и asyncio. Для базы Python + Redis выделяются Elasticsearch и Kafka. Отрицательная разница наблюдается у групп с HTML, CSS и Java, а для троек --- также при добавлении CSS к Python + SQL.

На результат влияет состав вакансий. В группе Bottle + Python все 11 вакансий с зарплатой требуют от трёх лет опыта. В группе Git + HTML 16 из 20 рассчитаны на опыт до трёх лет или без опыта. Поэтому разницу медиан нельзя приписать только навыку. Малые группы чувствительны к отдельным предложениям. Проверка состава групп приведена в \file{sql/experience.sql}.

\newpage\heading{6. Сравнение регионов}

\begin{Verbatim}[fontsize=\footnotesize]
python scripts/analysis.py regions
\end{Verbatim}


Выбраны три региона с наибольшим числом вакансий варианта № 2. Общие показатели рассчитаны по всем вакансиям региона, включая записи без навыков. Зарплатные показатели используют те же рублёвые суммы на руки.

\begingroup\small\setlength{\tabcolsep}{3pt}\renewcommand{\arraystretch}{1.18}
\begin{tabularx}{\linewidth}{Yrrrrr}
\toprule
Регион & Вакансий & С навыками & Зарплат & Доля, \% & Медиана, руб. \\
\midrule
Москва & 1\,908 & 1\,430 & 381 & 19,97 & 200\,000 \\
Санкт-Петербург & 525 & 403 & 144 & 27,43 & 150\,000 \\
Екатеринбург & 152 & 122 & 45 & 29,61 & 150\,000 \\
\bottomrule
\end{tabularx}
\endgroup


Доля --- часть всех вакансий региона с подходящей зарплатой. Для каждого размера ниже приведено самое частое сочетание. Частота считается по всем вакансиям региона, медиана --- по записям с зарплатой.

\begingroup\small\setlength{\tabcolsep}{3pt}\renewcommand{\arraystretch}{1.18}
\begin{tabularx}{\linewidth}{L{31mm}cYrrr}
\toprule
Регион & $k$ & Сочетание & Частота & Зарплат & Медиана, руб. \\
\midrule
Москва & 1 & python & 1359 & 302 & 200\,000 \\
Москва & 2 & postgresql + python & 613 & 139 & 200\,000 \\
Москва & 3 & docker + postgresql + python & 227 & 52 & 225\,000 \\
Санкт-Петербург & 1 & python & 391 & 111 & 140\,000 \\
Санкт-Петербург & 2 & postgresql + python & 190 & 55 & 170\,000 \\
Санкт-Петербург & 3 & docker + postgresql + python & 69 & 17 & 135\,000 \\
Екатеринбург & 1 & python & 114 & 31 & 150\,000 \\
Екатеринбург & 2 & postgresql + python & 55 & 15 & 150\,000 \\
Екатеринбург & 3 & git + postgresql + python & 22 & 5 & 108\,000 \\
\bottomrule
\end{tabularx}
\endgroup


\shot{11_regions}{Региональные показатели и частые сочетания}


Москва даёт наибольшую выборку и медиану 200\,000 руб. В Санкт-Петербурге и Екатеринбурге медиана равна 150\,000 руб., на 25\% ниже московской. Доля вакансий с подходящей зарплатой в Москве ниже: 19,97\% против 27,43\% и 29,61\%.

Во всех трёх регионах самая частая пара --- PostgreSQL + Python. Её доля --- 32,13\% в Москве и около 36,2\% в двух других регионах. Самая частая тройка в Москве и Санкт-Петербурге --- Docker + PostgreSQL + Python. В Екатеринбурге лидирует Git + PostgreSQL + Python, но её медиана рассчитана только по пяти зарплатам.

Для тройки Docker + PostgreSQL + Python медианы Москвы и Санкт-Петербурга равны 225\,000 и 135\,000 руб. по 52 и 17 зарплатам. Одинаковый стек не означает одинаковую медиану в разных регионах.

\block{7. Вывод}

Для варианта № 2 обработаны 3\,892 вакансии Python-разработчиков. MapReduce из S3 и независимый SQL-расчёт DuckDB полностью совпали. Повтор ошибочной задачи не изменил результат. DuckDB оказался быстрее по медиане времени, MapReduce потребовал меньше памяти по измеренной сумме RSS.

Зарплатные различия зависят от сочетания навыков, региона и состава вакансий. Для всех трёх регионов выделяется пара PostgreSQL + Python. Наибольшая общая медиана среди выбранных регионов получена для Москвы.

\end{document}
